\originalchapter{弹性、序参量与对称破缺}
\label{ch:landau}
\source{8.334早期讲义对应L1--L3；符号与量纲在此统一。}
\originalsection{从粒子到连续场}
晶体平衡位置$R_i$附近的位移$u_i$在长波可插值为$u(x)$。均匀平移$u\to u+a$不应改变能量，因此最低项含梯度。刚体转动也不产生弹性能，真正的应变为$u_{ij}=(\partial_iu_j+\partial_ju_i)/2$而不是所有位移梯度。各向同性小应变自由能
\begin{equation}
F_{\rm el}=\int\dd^dx\,[\lambda(\nabla\cdot u)^2/2+\mu u_{ij}u_{ij}],
\end{equation}
$\mu>0$且体积模量满足正条件才稳定。Fourier分解成纵横投影$P_L=\hat q\hat q^T,P_T=I-P_L$，静态关联为$\avg{u_i(q)u_j(-q)}=\kb T[P_L/(\lambda+2\mu)+P_T/\mu]_{ij}/q^2$。零波矢整体平移须固定或单独积分，否则Gaussian配分出现零模无限体积。

关联$\int\dd^dq/q^2$在$d\le2$红外发散，长波位移可使二维晶体的通常平移长程序变为代数序，但局部应变含额外$q$因子，仍可有限。这说明破缺哪种对称与观察哪个场同等重要，不能把大位移涨落一概称晶体马上熔化。
\originalsection{Landau--Ginzburg泛函}
序参量$\phi(x)$代表局域平均密度差、磁化或其他慢变量。短程相互作用、缓慢空间变化、相关对称约束允许
\begin{equation}
\mathcal H[\phi]\equiv\frac F{\kb T}=\int\dd^dx\left[\frac K2(\nabla\phi)^2+\frac r2\phi^2+\frac u4\phi^4-h\phi\right].
\end{equation}
这里$\mathcal H$无量纲，$K,r,u,h$已吸收温度与场的单位。场积分必须有微观截止$a$，不是未定义的无穷多个独立点积分。对Ising$\phi\to-\phi$，$h=0$禁奇项；对无这种对称的液气临界点可移场并调参，使最低临界形状同标量模型。$u>0$稳定四阶截断；若$u<0$需六阶正项，可能一阶转变或三临界。

均匀鞍点$r\phi+u\phi^3=h$给平均场指数。空间Euler方程$-K\nabla^2\phi+r\phi+u\phi^3=h$也给非均匀构造。例如$r<0,h=0$的平面畴壁$\phi(x)=\phi_0\tanh[x/(\sqrt2\xi)]$，$\phi_0=\sqrt{-r/u},\xi=\sqrt{K/|r|}$，代回可核对二阶导与局域势平衡。界面张力$\sigma=\kb T\int K(\partial_x\phi)^2\dd x=2\sqrt2\kb T\sqrt K|r|^{3/2}/(3u)$，不是单纯两极小值的自由能差。
\originalsection{连续对称与Goldstone模}
$n$分量$O(n)$模型将$\phi^2$换成$\vct\phi^2$。$r<0$的均匀极小在半径$\phi_0$球面上，选$n$方向后写纵$\sigma$与$n-1$横$\pi$。二次能量为$\frac12K(\nabla\sigma)^2+|r|\sigma^2+\frac12K(\nabla\pi)^2$。纵模有质量平方$2|r|/K$，横模没有局域恢复项，因为均匀转向不改能量。这是Goldstone模的具体来源。

热横涨落$\avg{\pi^2}\propto(n-1)T\int\dd^dq/(Kq^2)$，$d\le2$会在无限尺度增长。对有限温、短程、连续内部对称的低维体系，这提示通常长程序被破坏；严谨Mermin--Wagner定理还要求其标准相互作用与对称假设。离散Ising无连续转向Goldstone模，二维可以有有限温序。外场、各向异性和长程力可给质量或改变红外幂，不能套同一禁止结论。
\originalsection{临界模式与对称允许项}
平移、旋转、内部对称以及守恒量决定允许的低能项；忽略一个对称允许且相关的项会得到错误普适类。例如有各向异性$\sum_i\phi_i^4$时并非严格$O(n)$，而有长程势$|q|^\sigma$时梯度核不再$q^2$。Landau泛函的作用是把微观细节组织为耦合常数，下一步RG才决定哪些细节在长距离留下。仅写出四次多项式并不证明它是某材料的完整有效理论。
